\documentclass[10pt,letterpaper]{article}
\usepackage[utf8]{inputenc}
\usepackage{geometry}
\usepackage{hyperref}
\usepackage{enumitem}
\usepackage{titlesec}
\usepackage{charter} % Premium professional font

% Page setup
\geometry{letterpaper, left=0.55in, right=0.55in, top=0.5in, bottom=0.5in}
\pdfgentounicode=1 % Ensure PDF is searchable (important for ATS scanners)

% Color variables (Graphite theme)
\hypersetup{
    colorlinks=true,
    linkcolor=black,
    urlcolor=blue,
    pdftitle={Quantitative Developer Resume Template},
    pdfauthor={QuantMentor}
}

% Style definitions
\titleformat{\section}{\large\bfseries\scshape}{}{0em}{}[\titlerule]
\titlespacing{\section}{0pt}{10pt}{6pt}
\setlist[itemize]{noitemsep, topsep=2pt, leftmargin=15pt}

\begin{document}
\pagestyle{empty}

% Contact Header
\begin{center}
    {\Huge \textbf{AMIT JHA}} \\ \medskip
    \small
    \href{mailto:jha.8@alumni.iitj.ac.in}{jha.8@alumni.iitj.ac.in} \ | \ 
    +91 98765 43210 \ | \ 
    Mumbai, India \ | \ 
    \href{https://github.com/iitakjha97}{github.com/iitakjha97} \ | \ 
    \href{https://linkedin.com/in/amit-jha-quant}{linkedin.com/in/amit-jha-quant}
\end{center}

% Education Section
\section{Education}
\textbf{Indian Institute of Technology (IIT), Jodhpur} \hfill Jodhpur, India \\
\textit{Bachelor of Technology in Computer Science \& Engineering; GPA: 9.1/10.0} \hfill July 2019 -- May 2023
\begin{itemize}
    \item \textbf{Coursework}: Linear Algebra, Probability \& Statistics, Numerical Methods, Data Structures \& Algorithms.
    \item \textbf{Awards}: Institute Merit-cum-Means Scholarship recipient (Top 5\% of CSE department).
\end{itemize}

\medskip

\textbf{World Quantitative Academy (WQA)} \hfill Remote \\
\textit{Advanced Certificate in Quantitative Finance \& Volatility Modeling} \hfill June 2023 -- Nov 2023
\begin{itemize}
    \item Intensive coursework in Stochastic Calculus, Black-Scholes PDE discretization, and Monte Carlo methods.
\end{itemize}

% Professional Experience Section
\section{Professional Experience}
\textbf{Model Risk \& Quant Solutions Group} \hfill Mumbai, India \\
\textit{Quantitative Analyst / Model Validator} \hfill June 2023 -- Present
\begin{itemize}
    \item Perform rigorous mathematical validation on sell-side front-office pricing engines (Equity, FX, and Interest Rate).
    \item Validated Crank-Nicolson finite difference schemes for barrier options, verifying stability bounds under Rannacher time-stepping.
    \item Audited stochastic volatility models (Heston, SABR) by verifying asymptotic wing behaviors and local vol Dupire mappings.
    \item Developed independent Python libraries to replicate Monte Carlo paths using the Quadratic Exponential (QE) scheme.
\end{itemize}

\medskip

\textbf{High Frequency Trading Lab, IIT} \hfill Jodhpur, India \\
\textit{Quantitative Research Intern} \hfill May 2022 -- July 2022
\begin{itemize}
    \item Designed and backtested intraday statistical arbitrage strategies on NSE equities using order book imbalance features.
    \item Optimized a C++17 trading simulation engine using vectorization (Eigen library), reducing simulation latency by 35\%.
    \item Evaluated cointegration relationships (Johansen test) and calibrated mean-reversion rates using Ornstein-Uhlenbeck processes.
\end{itemize}

% Selected Projects Section
\section{Selected Projects}
\textbf{High-Performance Equity Options Pricing Engine (C++17)} \hfill \href{https://github.com/iitakjha97/quant-pde-solver}{\textit{GitHub Link}}
\begin{itemize}
    \item Developed a modular PDE solver using Finite Difference Methods (Explicit, Implicit, and Crank-Nicolson schemes).
    \item Implemented SOR (Successive Over-Relaxation) and Thomas algorithms in C++ for tridiagonal matrix inversion.
    \item Wrote custom Python bindings using pybind11 to visualize options Greeks surface grids in Jupyter notebooks.
\end{itemize}

\medskip

\textbf{Heston Model Volatility Calibration (Python, SciPy, PyTorch)} \hfill \href{https://github.com/iitakjha97/heston-calibration}{\textit{GitHub Link}}
\begin{itemize}
    \item Calibrated Heston model parameters ($v_0, \kappa, \theta, \sigma, \rho$) to market-implied volatility surfaces.
    \item Implemented the Carr-Madan FFT (Fast Fourier Transform) option pricing methodology using characteristic functions.
    \item Structured a deep PINN (Physics-Informed Neural Network) solver to compute option values directly from the BS PDE.
\end{itemize}

% Technical Skills Section
\section{Technical Skills}
\begin{itemize}[leftmargin=10pt]
    \item \textbf{Languages}: C++ (17/20, STL, Template Programming), Python (NumPy, SciPy, Pandas, PyTorch), SQL, Bash, LaTeX.
    \item \textbf{Quantitative Finance}: Black-Scholes PDE, Stochastic Volatility (SABR, Heston), Finite Differences, Monte Carlo.
    \item \textbf{Libraries \& Tools}: Eigen, Boost.Asio, Git, Docker, Linux (Ubuntu/CentOS), Visual Studio Code, CMake.
\end{itemize}

\end{document}
